% !TEX program = xelatex
\documentclass[12pt]{ctexart}

\usepackage[a4paper,margin=2.3cm,top=2.6cm]{geometry}
\usepackage{amsmath,amssymb,amsthm}
\usepackage{enumitem}
\usepackage{array}
\usepackage{fancyhdr}
\usepackage{xcolor}

\setlist[enumerate]{itemsep=10pt,topsep=6pt,leftmargin=2.2em}

\pagestyle{fancy}
\fancyhf{}
\lhead{\small 数学分析（上）期末考试 · B 卷}
\rhead{\small 第 \thepage 页}
\cfoot{}
\renewcommand{\headrulewidth}{0.4pt}

\newcommand{\bigsection}[2]{%
  \vspace{6pt}%
  {\noindent\bfseries\large #1}\hfill{\normalsize（本大题共 #2）}\par
  \vspace{2pt}\hrule height0.4pt\vspace{8pt}}

\newcommand{\blank}[1][3cm]{\underline{\hspace{#1}}}

\newcommand{\optgrid}[4]{%
  \par\smallskip\noindent
  \begin{tabular}{@{}p{0.45\textwidth}p{0.45\textwidth}@{}}
  (A)~#1 & (B)~#2 \\[2pt]
  (C)~#3 & (D)~#4
  \end{tabular}}

\linespread{1.15}

\begin{document}

%==================== 封面信息 ====================
\begin{center}
  {\Large\bfseries 数学分析（上册）期末考试试卷（B 卷）}\\[4pt]
  {\small 教材：华东师范大学数学系《数学分析》（第 5 版）\;\;|\;\;
  范围：导数与微分、微分中值定理及其应用、实数的完备性、不定积分、定积分、
  定积分应用、反常积分；多元函数的极限与连续、多元函数微分学}\\[10pt]
\end{center}

\noindent
\begin{tabular}{@{}p{4.5cm}p{4.5cm}p{4.5cm}@{}}
姓名：\blank[3cm] & 学号：\blank[3cm] & 班级：\blank[3cm]
\end{tabular}

\vspace{6pt}
\noindent\textbf{满分 100 分\qquad 考试时间 120 分钟\qquad 闭卷}

\vspace{8pt}
\noindent
\renewcommand{\arraystretch}{1.4}
\begin{tabular}{|c|c|c|c|c|c|c|c|}
\hline
题号 & 一 & 二 & 三 & 四 & 五 & 六 & 总分 \\\hline
满分 & 15 & 10 & 15 & 36 & 12 & 12 & 100 \\\hline
得分 & & & & & & & \\\hline
\end{tabular}

\vspace{14pt}

%==================== 一、选择题 ====================
\bigsection{一、选择题}{5 小题，每题 3 分，共 15 分。每题只有一个正确选项}

\begin{enumerate}[label=\arabic*.]
\item 下列函数中，在 $x=0$ 处可导的是（\qquad）
  \optgrid{$|x|$}{$x|x|$}{$\sqrt{|x|}$}{$\sqrt[3]{x}$}

\item 函数 $f(x)=\ln x$ 在 $[1,e]$ 上满足拉格朗日中值定理，其中值点 $\xi=$（\qquad）
  \optgrid{$1$}{$e-1$}{$\dfrac{e+1}{2}$}{$\sqrt{e}$}

\item 反常积分 $\displaystyle\int_{0}^{1}\frac{1}{x^{p}}\,\mathrm{d}x$ 收敛的充要条件是（\qquad）
  \optgrid{$p>1$}{$p<1$}{$p\geq 1$}{$p>0$}

\item 数列 $\{a_{n}\}$ 收敛的充要条件是（\qquad）
  \optgrid{$\{a_{n}\}$ 单调}{$\{a_{n}\}$ 有界}{$\{a_{n}\}$ 是柯西列}{$\{a_{n}\}$ 有收敛子列}

\item 设二元函数 $z=f(x,y)$ 在点 $P_{0}$ 处两个偏导数都存在，则在 $P_{0}$ 处（\qquad）
  \optgrid{$f$ 必连续}{$f$ 必可微}{$f$ 不一定连续}{$f$ 必取得极值}
\end{enumerate}

%==================== 二、判断题 ====================
\bigsection{二、判断题}{5 小题，每题 2 分，共 10 分。正确打 “$\checkmark$”，错误打 “$\times$”}

\begin{enumerate}[label=\arabic*.]
\item 若 $f$ 在点 $x_{0}$ 连续，则 $f$ 在点 $x_{0}$ 必可导。\hfill（\qquad）
\item 若 $f$ 在 $[a,b]$ 上可积，则 $f$ 在 $[a,b]$ 上必有界。\hfill（\qquad）
\item 若 $f$ 在 $[a,b]$ 上连续，则存在 $\xi\in[a,b]$ 使
  $\displaystyle\int_{a}^{b}f(x)\,\mathrm{d}x=f(\xi)(b-a)$。\hfill（\qquad）
\item 反常积分 $\displaystyle\int_{1}^{+\infty}\frac{1}{x}\,\mathrm{d}x$ 收敛。\hfill（\qquad）
\item 若二元函数在某点可微，则它在该点的两个偏导数都存在。\hfill（\qquad）
\end{enumerate}

%==================== 三、填空题 ====================
\bigsection{三、填空题}{5 小题，每题 3 分，共 15 分}

\begin{enumerate}[label=\arabic*.]
\item 设 $f(x)=\ln(1+x^{2})$，则 $f'(x)=\blank[3.5cm]$。
\item $\displaystyle\int \cos x\,\mathrm{d}x=\blank[4cm]$。
\item $\displaystyle\int_{0}^{\pi}\sin x\,\mathrm{d}x=\blank[3cm]$。
\item 反常积分 $\displaystyle\int_{1}^{+\infty}\frac{1}{x^{2}}\,\mathrm{d}x=\blank[3cm]$。
\item 设 $z=x^{2}+xy+y^{2}$，则 $\dfrac{\partial z}{\partial x}=\blank[3.5cm]$。
\end{enumerate}

%==================== 四、计算题 ====================
\bigsection{四、计算题}{6 小题，每题 6 分，共 36 分。要求写出主要计算步骤}

\begin{enumerate}[label=\arabic*.]
\item 求极限 $\displaystyle\lim_{x\to 0}\frac{1-\cos x}{x^{2}}$。

\item 求不定积分 $\displaystyle\int \ln x\,\mathrm{d}x$。

\item 计算定积分 $\displaystyle\int_{0}^{1} x\,e^{x^{2}}\,\mathrm{d}x$。

\item 求函数 $f(x)=\dfrac{x}{1+x^{2}}$ 的单调区间与极值。

\item 设 $z=x^{3}+y^{3}-3xy$，求一阶偏导数 $z_{x},z_{y}$ 及全部二阶偏导数
  $z_{xx},z_{yy},z_{xy}$。

\item 计算反常积分 $\displaystyle\int_{0}^{1}\frac{1}{\sqrt{1-x^{2}}}\,\mathrm{d}x$。
\end{enumerate}

%==================== 五、证明题 ====================
\bigsection{五、证明题}{2 小题，每题 6 分，共 12 分}

\begin{enumerate}[label=\arabic*.]
\item 证明：当 $x>0$ 时，$e^{x}>1+x$。

\item 设 $a_{1}=\sqrt{2}$，$a_{n+1}=\sqrt{2+a_{n}}\ (n=1,2,\dots)$。
  利用单调有界定理证明数列 $\{a_{n}\}$ 收敛，并求其极限。
\end{enumerate}

%==================== 六、应用题 ====================
\bigsection{六、应用题}{2 小题，每题 6 分，共 12 分}

\begin{enumerate}[label=\arabic*.]
\item 求由曲线 $y=\sqrt{x}$、直线 $x=4$ 及 $x$ 轴所围平面图形的面积 $A$，
  并求该图形绕 $x$ 轴旋转一周所得旋转体的体积 $V$。

\item 用长为 $20\ \mathrm{m}$ 的篱笆，靠一面墙围成一块矩形场地（靠墙一侧不用篱笆）。
  问矩形的长、宽各为多少时围成的面积最大？最大面积是多少？
\end{enumerate}

\vspace{8pt}
\begin{center}\small —— 试卷结束，参考答案见后 ——\end{center}

%==================== 参 考 答 案 ====================
\clearpage
\begin{center}
  {\Large\bfseries 参考答案与评分标准（B 卷）}\\[2pt]
  {\small 数学分析（上册）期末考试}
\end{center}
\vspace{6pt}\hrule\vspace{12pt}

\noindent\textbf{一、选择题}（每题 3 分）\par\vspace{4pt}
\begin{center}
\renewcommand{\arraystretch}{1.3}
\begin{tabular}{|c|c|c|c|c|c|}
\hline
题号 & 1 & 2 & 3 & 4 & 5 \\\hline
答案 & B & B & B & C & C \\\hline
\end{tabular}
\end{center}
\noindent\small 简析：1.\,$x|x|$ 在 $0$ 处左右导数均为 $0$，可导；其余三者在 $0$ 处不可导。
2.\,$\dfrac{1}{\xi}=\dfrac{\ln e-\ln 1}{e-1}=\dfrac{1}{e-1}\Rightarrow\xi=e-1\in(1,e)$。
3.\,瑕点在 $x=0$，$p$-积分当且仅当 $p<1$ 收敛。
4.\,实数系中数列收敛 $\Leftrightarrow$ 为柯西列（柯西收敛准则）。
5.\,偏导数存在不能保证连续（更不能保证可微）。\normalsize

\vspace{10pt}
\noindent\textbf{二、判断题}（每题 2 分）\par\vspace{4pt}
\begin{center}
\renewcommand{\arraystretch}{1.3}
\begin{tabular}{|c|c|c|c|c|c|}
\hline
题号 & 1 & 2 & 3 & 4 & 5 \\\hline
答案 & $\times$ & $\checkmark$ & $\checkmark$ & $\times$ & $\checkmark$ \\\hline
\end{tabular}
\end{center}
\noindent\small 简析：1.\,连续未必可导，反例 $|x|$ 在 $0$ 处。
2.\,可积的必要条件是有界。3.\,积分第一中值定理。
4.\,$\int_{1}^{+\infty}\frac1x\mathrm{d}x$ 发散（$p=1$）。
5.\,可微 $\Rightarrow$ 偏导数存在。\normalsize

\vspace{10pt}
\noindent\textbf{三、填空题}（每题 3 分）\par\vspace{4pt}
\begin{enumerate}[label=\arabic*.,itemsep=4pt]
\item $\dfrac{2x}{1+x^{2}}$ \qquad
\item $\sin x+C$ \qquad
\item $2$（$[-\cos x]_{0}^{\pi}=1+1=2$）
\item $1$（$[-x^{-1}]_{1}^{+\infty}=0-(-1)=1$）
\item $2x+y$
\end{enumerate}

\vspace{6pt}
\noindent\textbf{四、计算题}（每题 6 分）

\begin{enumerate}[label=\textbf{\arabic*.},itemsep=10pt]
\item 由 $1-\cos x=\dfrac{x^{2}}{2}+o(x^{2})$，得
  $\displaystyle\lim_{x\to0}\frac{1-\cos x}{x^{2}}=\frac12$。
  （或用洛必达：$\lim\dfrac{\sin x}{2x}=\dfrac12$。）

\item 用分部积分，取 $u=\ln x,\ \mathrm{d}v=\mathrm{d}x$，则 $v=x$：
  \[
  \int\ln x\,\mathrm{d}x=x\ln x-\int x\cdot\frac1x\,\mathrm{d}x=x\ln x-x+C .
  \]

\item 令 $u=x^{2}$，$\mathrm{d}u=2x\,\mathrm{d}x$：
  \[
  \int_{0}^{1}x\,e^{x^{2}}\,\mathrm{d}x=\frac12\int_{0}^{1}e^{u}\,\mathrm{d}u
  =\frac12\big(e^{u}\big)\Big|_{0}^{1}=\frac{e-1}{2}.
  \]

\item $f'(x)=\dfrac{(1+x^{2})-x\cdot 2x}{(1+x^{2})^{2}}=\dfrac{1-x^{2}}{(1+x^{2})^{2}}$。
  令 $f'(x)=0$ 得 $x=\pm 1$。
  \[
  \begin{array}{c|ccccc}
   x & (-\infty,-1) & -1 & (-1,1) & 1 & (1,+\infty)\\\hline
   f'(x) & - & 0 & + & 0 & -\\
   f(x) & \searrow & \text{极小} & \nearrow & \text{极大} & \searrow
  \end{array}
  \]
  单调增区间 $(-1,1)$，单调减区间 $(-\infty,-1)$ 与 $(1,+\infty)$；
  极小值 $f(-1)=-\dfrac12$，极大值 $f(1)=\dfrac12$。

\item 一阶偏导：$z_{x}=3x^{2}-3y,\qquad z_{y}=3y^{2}-3x$。\\
  二阶偏导：
  \[
  z_{xx}=6x,\quad z_{yy}=6y,\quad z_{xy}=z_{yx}=-3.
  \]

\item 瑕点在 $x=1$：
  \[
  \int_{0}^{1}\frac{\mathrm{d}x}{\sqrt{1-x^{2}}}
  =\lim_{b\to 1^{-}}\big[\arcsin x\big]_{0}^{b}
  =\lim_{b\to1^{-}}\arcsin b=\frac{\pi}{2}.
  \]
\end{enumerate}

\vspace{6pt}
\noindent\textbf{五、证明题}（每题 6 分）

\begin{enumerate}[label=\textbf{\arabic*.},itemsep=12pt]
\item 令 $g(x)=e^{x}-1-x$，则 $g(0)=0$，且当 $x>0$ 时
  $g'(x)=e^{x}-1>0$，故 $g$ 在 $[0,+\infty)$ 上单调递增。
  于是当 $x>0$ 时 $g(x)>g(0)=0$，即 $e^{x}>1+x$。$\qquad\blacksquare$

\item \textbf{有界：} 用归纳法证 $0<a_{n}<2$。$a_{1}=\sqrt2<2$；
  设 $a_{n}<2$，则 $a_{n+1}=\sqrt{2+a_{n}}<\sqrt{2+2}=2$。故对一切 $n$ 有 $0<a_{n}<2$，
  即 $\{a_{n}\}$ 有上界。\\[2pt]
  \textbf{单调递增：} 由
  $a_{n+1}^{2}-a_{n}^{2}=(2+a_{n})-a_{n}^{2}=-(a_{n}-2)(a_{n}+1)$，
  而 $0<a_{n}<2$ 时 $(a_{n}-2)<0,\ (a_{n}+1)>0$，故 $a_{n+1}^{2}-a_{n}^{2}>0$；
  又 $a_{n},a_{n+1}>0$，所以 $a_{n+1}>a_{n}$，数列单调递增。\\[2pt]
  \textbf{求极限：} 由\textbf{单调有界定理}，$\{a_{n}\}$ 收敛，设
  $\displaystyle\lim_{n\to\infty}a_{n}=L$。在 $a_{n+1}=\sqrt{2+a_{n}}$ 两边取极限得
  $L=\sqrt{2+L}$，即 $L^{2}-L-2=0$，$(L-2)(L+1)=0$。
  因 $L>0$，故 $L=2$，即 $\displaystyle\lim_{n\to\infty}a_{n}=2$。$\qquad\blacksquare$
\end{enumerate}

\vspace{6pt}
\noindent\textbf{六、应用题}（每题 6 分）

\begin{enumerate}[label=\textbf{\arabic*.},itemsep=12pt]
\item 在 $[0,4]$ 上 $y=\sqrt{x}\geq 0$，面积
  \[
  A=\int_{0}^{4}\sqrt{x}\,\mathrm{d}x=\Big[\frac23 x^{3/2}\Big]_{0}^{4}
  =\frac23\cdot 8=\frac{16}{3}.
  \]
  绕 $x$ 轴旋转（圆盘法）：
  \[
  V=\pi\int_{0}^{4}(\sqrt{x})^{2}\,\mathrm{d}x=\pi\int_{0}^{4}x\,\mathrm{d}x
  =\pi\Big[\frac{x^{2}}{2}\Big]_{0}^{4}=8\pi.
  \]

\item 设垂直于墙的两条边长为 $x$（$0<x<10$），则平行于墙的一边长为 $20-2x$，
  面积
  \[
  S(x)=x(20-2x)=20x-2x^{2}.
  \]
  $S'(x)=20-4x=0\Rightarrow x=5$；当 $x<5$ 时 $S'>0$，$x>5$ 时 $S'<0$，
  故 $x=5$ 为最大值点。此时平行墙的一边长 $20-2\cdot5=10$，最大面积
  \[
  S(5)=5\times 10=50\ (\mathrm{m}^{2}).
  \]
  即矩形宽（垂直墙）$5\,\mathrm{m}$、长（平行墙）$10\,\mathrm{m}$ 时面积最大，
  为 $50\,\mathrm{m}^{2}$。
\end{enumerate}

\end{document}
